% !PRINTED 388
\weberpage{388}
\begin{center}
§ 106. \gesperrt{Die Geschlechter der Idealklassen.}
\end{center}

In \S\,100 haben wir gesehen, wie wir aus einer Zahlgruppe
im Körper $\Omega$ Einteilungen der Ideale in Klassen ableiten können,
und wie für diese Einteilungen die Klassenzahl zu bestimmen ist.

Wir wollen hier auch die Einteilung in \gesperrt{Geschlechter} aus
diesem Gesichtspunkte betrachten.

Um unsere Betrachtungen gleich auf die Ordnungen ausdehnen
 zu können, scheiden wir von dem Gebiet der \gesperrt{natürlichen
 Zahlen} zunächst alle Zahlen aus, die mit irgend einer
beliebig angenommenen Zahl $S$ einen gemeinschaftlichen Teiler
haben.

Hierauf nehmen wir einen positiven rationalen Modul $m$, und
nehmen in $S$ unter anderen alle Primzahlen auf, die in $m$ enthalten
 sind.

Damit sind also alle Zahlen ausgeschieden, die einen gemeinschaftlichen
 Teiler mit $m$ haben.

Aus den übrigbleibenden Zahlen und ihren Quotienten
kann man eine \gesperrt{Gruppe rationaler Brüche} bilden, die wir mit
$Z$ bezeichnen wollen.

\gesperrt{Ersetzen wir (nach Gauss) eine Zahl $a^{-1}$ durch die
ganze Zahl, die, mit $a$ multipliziert, eine der Einheit
kongruente Zahl ergibt, so können wir die Lehre von
der Kongruenz nach dem Modul $m$ ohne weiteres auf
die gebrochenen Zahlen $Z$ übertragen.}

Sind dann $a/b$ und $a_1/b_1$ Brüche in $Z$, so ist $a/b \equiv a_1/b_1$,
wenn $ab_1 \equiv a_1 b$ im gewöhnlichen Sinne ist.

Vereinigen wir also alle untereinander kongruenten Zahlen
in $Z$ in eine Klasse, so zerfällt $Z$ in $\varphi(m)$-Zahlklassen, deren
jede durch eine ganze rationale Zahl repräsentiert werden kann.

Ist $M$ die Gruppe der Zahlen $a$ aus $Z$, die der Bedingung
\[
a \equiv 1 \pmod m
\]
genügen, so sind die Zahlklassen (modulo $m$) die Nebengruppen
zu $M$, und es ist
\weq{1}{(Z, M) = \varphi(m).}
Wir betrachten nun, wie im \S\,98, die Gruppen $O$ der ganzen
und gebrochenen Zahlen $\omega$ des quadratischen Körpers $\Omega$, \gesperrt{oder
einer Ordnung $[Q]$ dieses Körpers, von denen wir alle}
%CONT>>
